\capitulocapa{images/cap-8.png}{Como atualizar estado sem espalhar regras?}
\label{chap:acoes-atualizacoes}

\section{O filtro que cada tela atualizava de um jeito}

\begin{historia}
O Dashboard do DevPanel ganhou uma central de filtros. A primeira versão tinha
três controles: busca por texto, situação da conta e período de atividade.

A equipe criou uma store, expôs o objeto \codigo{filters} e uma ação genérica
chamada \codigo{setFilters}. O painel lateral usava essa ação para trocar a
situacao, a busca montava outro objeto e o botão ``Limpar'' decidia quais valores
deveriam voltar ao padrão.

Tudo funcionou até o período padrão mudar de sete para trinta dias. O painel foi
atualizado, mas o botão de limpar continuou restaurando sete. Uma exportação
rápida ainda apagava o período por completo. A store era a fonte única dos
dados, mas as regras estavam duplicadas nos consumidores.
\end{historia}

\imagemcapitulo
  {ch08-fluxo-acao-dominio-review}
  {Ações de domínio substituem cópias divergentes por um único próximo estado.}
  {fluxo-acao-dominio}

Mover um valor para Zustand resolve onde ele vive. Ainda precisamos decidir
quem conhece as regras de mudança. Se cada componente constrói o próximo estado,
a store vira apenas um recipiente compartilhado.

\section{O atalho habitual: um setter genérico}

É comum começar com um contrato semelhante a este:

\begin{lstlisting}[caption={Um contrato que entrega as regras aos consumidores.}]
type DashboardFilterStore = {
  filters: DashboardFilters
  setFilters: (filters: DashboardFilters) => void
}
\end{lstlisting}

O componente precisa conhecer o objeto inteiro para alterar uma parte:

\begin{lstlisting}
setFilters({
  ...filters,
  status: 'active',
})
\end{lstlisting}

O código parece pequeno, mas transfere duas responsabilidades para a tela. Ela
precisa saber como copiar o estado atual e quais combinações são válidas. Outro
componente pode esquecer o espalhamento e apagar campos que não pretendia mudar.

Uma variação expõe \codigo{setStatus}, \codigo{setSearch} e
\codigo{setPeriod}. Os nomes são melhores, porém ainda podem representar somente
atribuições. O passo importante não é trocar ``set'' por outro prefixo; é fazer
a ação expressar uma intenção e proteger uma regra.

\section{Definindo o estado antes das ações}

Os filtros do Dashboard possuem um vocabulário pequeno:

\begin{lstlisting}[caption={Tipos e valores iniciais dos filtros.},label={lst:dashboard-filter-types}]
type AccountStatus = 'all' | 'active' | 'invited' | 'blocked'
type ActivityPeriod = '7d' | '30d' | '90d'

type DashboardFilters = {
  search: string
  status: AccountStatus
  period: ActivityPeriod
}

const DEFAULT_FILTERS: DashboardFilters = {
  search: '',
  status: 'all',
  period: '30d',
}
\end{lstlisting}

O período padrão aparece uma única vez. Esse objeto descreve o estado para o qual
o produto deseja voltar, não apenas os valores usados na primeira renderização.

Poderíamos guardar cada campo no primeiro nível da store. Agrupá-los em
\codigo{filters} comunica que formam uma unidade de domínio e permite entregar
esse objeto à camada que monta a consulta. A escolha traz uma consequência:
como \codigo{filters} é um objeto aninhado, suas atualizações precisam criar um
novo objeto.

\section{Ações que preservam o contrato}

Agora definimos as operações permitidas:

\begin{lstlisting}[caption={Contrato orientado às intenções do Dashboard.}]
type DashboardFilterActions = {
  changeSearch: (search: string) => void
  selectStatus: (status: AccountStatus) => void
  selectPeriod: (period: ActivityPeriod) => void
  clearFilters: () => void
}

type DashboardFilterStore = {
  filters: DashboardFilters
} & DashboardFilterActions
\end{lstlisting}

Essas ações não expõem a possibilidade de substituir o objeto inteiro por uma
combinação arbitrária. Cada uma recebe apenas a informação necessária para a
interação que representa.

O contrato também ajuda o TypeScript. \codigo{selectPeriod('today')} não
compila, pois ``today'' não pertence ao conjunto de períodos aceitos pelo
produto. Ainda assim, tipos não substituem regras: é \codigo{clearFilters} que
decide como o estado deve ser restaurado.

\begin{decisao}
Exponha operações que os consumidores podem solicitar com segurança. Quando uma
regra mudar, deve existir um lugar óbvio para atualizá-la.
\end{decisao}

\section{Implementando atualizações funcionais}

A store completa usa a forma funcional de \codigo{set} para produzir cada novo
objeto a partir da fotografia mais recente:

\begin{lstlisting}[caption={Store dos filtros com ações semânticas.},label={lst:dashboard-filter-store}]
import { create } from 'zustand'

export const useDashboardFilterStore =
  create<DashboardFilterStore>()((set) => ({
    filters: DEFAULT_FILTERS,

    changeSearch: (search) =>
      set((state) => ({
        filters: { ...state.filters, search },
      })),

    selectStatus: (status) =>
      set((state) => ({
        filters: { ...state.filters, status },
      })),

    selectPeriod: (period) =>
      set((state) => ({
        filters: { ...state.filters, period },
      })),

    clearFilters: () =>
      set({ filters: { ...DEFAULT_FILTERS } }),
  }))
\end{lstlisting}

Em \codigo{changeSearch}, o próximo valor depende do restante do objeto atual.
A função recebida por \codigo{set} acessa a versão mais recente da store, copia
\codigo{state.filters} e substitui apenas \codigo{search}.

Não precisamos espalhar toda a store. Zustand combina por padrão a parcela
devolvida por \codigo{set} no primeiro nível. Essa combinação é rasa: ao trocar
\codigo{filters}, precisamos copiar explicitamente os campos internos que devem
permanecer.

\codigo{clearFilters} não depende dos filtros atuais. Ela cria um novo objeto a
partir de \codigo{DEFAULT\_FILTERS}. Assim, busca e situação são limpas e o
período volta para trinta dias. Quando o produto alterar o padrão novamente, a
inicialização e a limpeza continuarão alinhadas.

\imagemcapitulo
  {ch08-objeto-imutabilidade-review}
  {Uma atualização imutável preserva o objeto anterior e cria uma nova referência.}
  {objeto-imutabilidade}

\section{Por que não alterar o objeto diretamente?}

Objetos JavaScript são mutáveis, então este código é possível, mas está errado:

\begin{lstlisting}
selectStatus: (status) =>
  set((state) => {
    state.filters.status = status
    return { filters: state.filters }
  })
\end{lstlisting}

O código modifica a fotografia atual e devolve a mesma referência de
\codigo{filters}. Isso torna mudanças mais difíceis de rastrear e pode impedir
que assinaturas baseadas em igualdade reconheçam um novo valor.

Imutabilidade aqui não significa copiar a aplicação inteira. Copiamos apenas o
caminho que muda. Para um objeto com um nível, o espalhamento é suficiente. Se
o modelo exigir várias camadas de cópia em toda ação, vale primeiro perguntar se
o estado está aninhado demais antes de adicionar outra ferramenta.

\section{Conectando a central de filtros}

O painel visual assina os campos e ações que utiliza:

\begin{lstlisting}[caption={Consumidor envia intenções, sem montar o próximo estado.}]
export function FiltersPanel() {
  const filters = useDashboardFilterStore(
    (store) => store.filters,
  )
  const selectStatus = useDashboardFilterStore(
    (store) => store.selectStatus,
  )
  const selectPeriod = useDashboardFilterStore(
    (store) => store.selectPeriod,
  )
  const clearFilters = useDashboardFilterStore(
    (store) => store.clearFilters,
  )

  return (
    <aside>
      <StatusFilter
        value={filters.status}
        onChange={selectStatus}
      />
      <PeriodFilter
        value={filters.period}
        onChange={selectPeriod}
      />
      <button type="button" onClick={clearFilters}>
        Limpar filtros
      </button>
    </aside>
  )
}
\end{lstlisting}

O componente conhece os controles e os eventos da interface. Ele não sabe qual
é o período padrão nem como preservar outros campos. Se a regra de limpeza
mudar, esse JSX permanece igual.

Selecionar o objeto completo é aceitável neste momento porque o painel exibe
todos os filtros. Componentes menores devem selecionar somente o valor que
usam. No próximo capítulo veremos com cuidado como a forma de um selector afeta
renderizações.

\section{Valores derivados não precisam de uma ação}

O botão pode informar quantos filtros estão realmente ativos. Seria tentador
guardar \codigo{activeFilterCount} na store e atualizá-lo em todas as ações. Isso
criaria duas fontes para a mesma verdade.

O total pode ser calculado a partir dos filtros:

\begin{lstlisting}[caption={Valor derivado calculado a partir da fonte de verdade.}]
export function getActiveFilterCount(filters: DashboardFilters) {
  return Number(filters.search.trim() !== '')
    + Number(filters.status !== DEFAULT_FILTERS.status)
    + Number(filters.period !== DEFAULT_FILTERS.period)
}

const activeFilterCount = useDashboardFilterStore(
  (store) => getActiveFilterCount(store.filters),
)
\end{lstlisting}

Não existe ação \codigo{setActiveFilterCount}. Qualquer mudança em
\codigo{filters} produz o total correspondente. A regra permanece testável como
uma função simples e não corre o risco de ficar dessincronizada.

Nem todo cálculo deve ser colocado dentro do arquivo da store. Ele pertence
aqui porque traduz o significado de ``filtro ativo'' e é reutilizável. Uma
formatação usada por um único componente pode continuar próxima desse
componente.

\section{Até onde uma ação deve ir?}

Centralizar regras não significa esconder toda a aplicação dentro da store. A
ação deve coordenar mudanças do estado que pertence àquele domínio.

\codigo{clearFilters} pode restaurar os filtros. Ela não precisa manipular o
DOM, navegar para outra rota nem fazer uma consulta diretamente. A camada de
consulta pode observar os filtros e buscar os dados correspondentes; o botão
pode cuidar de foco e acessibilidade.

Também devemos evitar uma ação única como \codigo{updateDashboard}. Um nome
amplo não revela o evento, aceita responsabilidades futuras demais e devolve ao
consumidor a obrigação de conhecer a estrutura interna.

Uma boa fronteira deixa respostas simples:

\begin{itemize}
  \item a interface traduz eventos em chamadas de ações;
  \item a store preserva as regras de mudança de seus dados;
  \item valores derivados são calculados da fonte de verdade;
  \item efeitos externos permanecem na camada apropriada.
\end{itemize}

\section{Exercício: limpe sem perder o padrão}

\begin{tarefa}
O produto pediu que o botão ``Limpar filtros'' remova a busca e volte a situação
para ``Todas'', mas mantenha o período padrão de trinta dias. Implemente
\codigo{clearFilters} sem chamar três ações em sequência e sem repetir os valores
padrão dentro do componente.
\end{tarefa}

Depois da implementação, registre:

\begin{enumerate}
  \item onde o período padrão foi definido;
  \item por que o componente não recebe um \codigo{setFilters} genérico;
  \item quais referências são novas depois da limpeza;
  \item por que o total de filtros ativos não foi armazenado;
  \item qual responsabilidade deve continuar fora da store.
\end{enumerate}

\espacoresposta{9}

\begin{resumo}
\begin{itemize}
  \item Uma fonte única de estado não garante uma fonte única de regras.
  \item Ações semânticas descrevem intenções e protegem invariantes.
  \item A forma funcional de \codigo{set} usa a fotografia mais recente.
  \item Zustand combina atualizações apenas no primeiro nível da store.
  \item Objetos aninhados devem ser atualizados com novas referências.
  \item Valores calculáveis devem ser derivados, não sincronizados manualmente.
  \item A store coordena seu domínio sem absorver interface e efeitos externos.
\end{itemize}
\end{resumo}

No próximo capítulo, os indicadores do Dashboard mostrarão outro detalhe: mesmo
com ações bem definidas, um selector amplo pode fazer componentes atualizarem
quando o valor que exibem não mudou.
